% GENERATED FILE. Edit canonical lessons, then run npm run generate:books.
% canonical-source-hash: fnv1a64:3ef640a7b61c2888
% canonical-lessons: ZH-C207-didian, ZH-W207-di-ground, ZH-C207-jihua, ZH-C207-anpai, ZH-C207-zhaoxiang, ZH-C207-jingdian

\chapter{Where and When}
\label{ch:zh-where-and-when}
\hlchaptermodality{207}

\begin{chapteropening}
\textbf{By the end of this chapter:} \emph{I can agree on a place to meet, talk about plans, arrange things, take photos and visit the sights.}

Complete the last lesson of chapter 207: I can agree on a place to meet, talk about plans, arrange things, take photos and visit the sights.
\end{chapteropening}

\section[\texorpdfstring{\zh{地点} (dìdiǎn)}{地点 (dìdiǎn)}]{\zh{地点} (dìdiǎn) — place; venue}

\label{lesson:ZH-C207-didian}



\begin{quote}
Before the new one: say the Mandarin for office, then the Mandarin for reception desk.
\end{quote}

\subsection*{You'll want to know: \zh{地点}}
\textbf{\zh{地点}} — \emph{dìdiǎn} — \textquotedblleft{}place; venue\textquotedblright{}.

\begin{quote}
\textbf{\zh{我们见面的地点离这儿两公里}} \emph{wǒmen jiànmiàn de dìdiǎn lí zhèr liǎng gōnglǐ} — The place where we're meeting is two kilometres from here.
\end{quote}

The new character, in pieces: \zh{地} — six strokes: on the left the earth piece, a short bar, a vertical and a stroke rising to the right, as at the lower right of \zh{在}; on the right \zh{也}. The next lesson puts it in your hand.

\begin{cousinweb}
\textbf{\zh{地点}} (\emph{dìdiǎn}) is built from \zh{地} 'ground, place' and \zh{点} 'a point', both inherited from Middle Chinese, \emph{dijH} and \emph{temX} (Guangyun, in Baxter's transcription), which glosses the first 'the earth'.
\end{cousinweb}

\begin{grammarlens}[title={What you've built}]
One more word for talking about plans.
\end{grammarlens}

\subsection*{Your turn}
\begin{itemize}
\raggedright
  \item \emph{Say it:} \emph{dìdiǎn}
  \item \emph{Say it:} \emph{dìdiǎn}, once more
  \item \emph{Say it:} \emph{wǒmen jiànmiàn de dìdiǎn lí zhèr liǎng gōnglǐ} — The place where we're meeting is two kilometres from here.
  \item \emph{Recall:} say the Mandarin for office, then the Mandarin for reception desk, then say \emph{dìdiǎn} again
\end{itemize}

\begin{tcolorbox}[breakable,colback=teal!4,colframe=teal!35!black,title={Before you move on}]
What does \zh{地点} mean? (\textquotedblleft{}Place; venue\textquotedblright{}.) Say it once more.
\end{tcolorbox}
\section[\texorpdfstring{\zh{地} (dì)}{地 (dì)}]{\zh{地} — the earth piece, then \zh{也}}

\label{lesson:ZH-W207-di-ground}



\begin{quote}
The word first, then the character inside it.

\begin{itemize}
\raggedright
  \item \emph{Recall:} say \emph{dìdiǎn} — \textbf{R1}, one lesson back, before it decays
\end{itemize}

\begin{quote}
\zh{也} \zh{在}
\end{quote}

\emph{Read it:} each one aloud, with nothing said first — \emph{yě}, \emph{zài}
\end{quote}

\subsection*{Script — the earth piece, then \zh{也}}
\hlblockfigure{\detokenize{figures/ZH-W207-di-ground-filmstrip.pdf}}{How \zh{地} is written, stroke by stroke}

\begin{quote}
\zh{地}
\end{quote}

\textbf{Two pieces, left before right.}

\textbf{The left piece, the earth piece, three strokes.} A short bar, a vertical down through it, then a stroke along the foot that rises to the right. It is the piece at the lower right of \zh{在}, squeezed to the left.

\textbf{The right piece, \zh{也}, three strokes,} the same three as in \zh{也} itself: a bar that turns down and hooks, a short vertical, and a long vertical that bends right along the bottom and hooks up.

3 + 3 = six.

\subsection*{Your turn}
\begin{itemize}
\raggedright
  \item \emph{Write it:} \zh{地}, all six, with the model in front of you
  \item \emph{Cover:} write it again from memory, and count the pieces, not the strokes
\end{itemize}

\begin{tcolorbox}[breakable,colback=teal!4,colframe=teal!35!black,title={Before you move on}]
How many strokes in \zh{地}? (\textbf{Six}.) Which stroke of the left piece rises? (\textbf{The last, along the foot}.) What is the right piece? (\textbf{\zh{也}}.)

Source: \href{https://raw.githubusercontent.com/chanind/hanzi-writer-data/68d10a4b21150cae5e1ebbd223eed289cf32d90c/data/\%E5\%9C\%B0.json}{Hanzi Writer Data \zh{地}, pinned stroke paths and medians}.
\end{tcolorbox}
\section[jìhuà]{jìhuà — plan}

\label{lesson:ZH-C207-jihua}



\begin{quote}
Before the new one: say the Mandarin for reception desk, then the Mandarin for venue.
\end{quote}

\subsection*{You'll want to know: jìhuà}
\emph{jìhuà} — \textquotedblleft{}plan\textquotedblright{}.

\begin{quote}
\emph{nǐ zhōumò yǒu shénme jìhuà} — What are your plans for the weekend?
\end{quote}

Its characters come later; for now it is a word for the ear.

\begin{cousinweb}
\emph{jìhuà} is built from \emph{jì} 'to reckon' and \emph{huà} 'to mark out', both inherited from Middle Chinese, \emph{kejH} and \emph{hweak} (Guangyun, in Baxter's transcription), which glosses the first 'to count up' and the second 'to cut with an awl'.
\end{cousinweb}

\begin{grammarlens}[title={What you've built}]
One more word for talking about plans.
\end{grammarlens}

\subsection*{Your turn}
\begin{itemize}
\raggedright
  \item \emph{Say it:} \emph{jìhuà}
  \item \emph{Say it:} \emph{jìhuà}, once more
  \item \emph{Say it:} \emph{nǐ zhōumò yǒu shénme jìhuà} — What are your plans for the weekend?
  \item \emph{Recall:} say the Mandarin for reception desk, then the Mandarin for venue, then say \emph{jìhuà} again
\end{itemize}

\begin{tcolorbox}[breakable,colback=teal!4,colframe=teal!35!black,title={Before you move on}]
What does jìhuà mean? (\textquotedblleft{}Plan\textquotedblright{}.) Say it once more.
\end{tcolorbox}
\section[ānpái]{ānpái — to arrange}

\label{lesson:ZH-C207-anpai}



\begin{quote}
Before the new one: say the Mandarin for venue, then the Mandarin for a plan.
\end{quote}

\subsection*{You'll want to know: ānpái}
\emph{ānpái} — \textquotedblleft{}to arrange\textquotedblright{}.

\begin{quote}
\emph{māma ānpáile míngtiān de shíjiān} — Mum has planned out tomorrow.
\end{quote}

Its characters come later; for now it is a word for the ear.

\begin{cousinweb}
\emph{ānpái} is built from \emph{ān} 'to settle' and \emph{pái} 'to set in a row'. The first is inherited from Middle Chinese \emph{'an} (Guangyun, in Baxter's transcription), glossed 'calm'; the dictionary glosses the second 'to push aside', so how the pair came to mean arranging is not settled.
\end{cousinweb}

\begin{grammarlens}[title={What you've built}]
One more word for talking about plans.
\end{grammarlens}

\subsection*{Your turn}
\begin{itemize}
\raggedright
  \item \emph{Say it:} \emph{ānpái}
  \item \emph{Say it:} \emph{ānpái}, once more
  \item \emph{Say it:} \emph{māma ānpáile míngtiān de shíjiān} — Mum has planned out tomorrow.
  \item \emph{Recall:} say the Mandarin for venue, then the Mandarin for a plan, then say \emph{ānpái} again
\end{itemize}

\begin{tcolorbox}[breakable,colback=teal!4,colframe=teal!35!black,title={Before you move on}]
What does ānpái mean? (\textquotedblleft{}To arrange\textquotedblright{}.) Say it once more.
\end{tcolorbox}
\section[zhàoxiàng]{zhàoxiàng — to take a photo}

\label{lesson:ZH-C207-zhaoxiang}



\begin{quote}
Before the new one: say the Mandarin for a plan, then the Mandarin for to arrange.
\end{quote}

\subsection*{You'll want to know: zhàoxiàng}
\emph{zhàoxiàng} — \textquotedblleft{}to take a photo\textquotedblright{}.

\begin{quote}
\emph{zhèr kěyǐ zhàoxiàng ma} — Can I take photos here?
\end{quote}

Its characters come later; for now it is a word for the ear.

\begin{cousinweb}
\emph{zhàoxiàng} is built from \emph{zhào} 'to shine on' and \emph{xiàng} 'a likeness', both inherited from Middle Chinese, \emph{tsyewH} and \emph{sjangH} (Guangyun, in Baxter's transcription), which glosses the first 'bright' and the second 'to look at'.
\end{cousinweb}

\begin{grammarlens}[title={What you've built}]
One more word for talking about plans.
\end{grammarlens}

\subsection*{Your turn}
\begin{itemize}
\raggedright
  \item \emph{Say it:} \emph{zhàoxiàng}
  \item \emph{Say it:} \emph{zhàoxiàng}, once more
  \item \emph{Say it:} \emph{zhèr kěyǐ zhàoxiàng ma} — Can I take photos here?
  \item \emph{Recall:} say the Mandarin for a plan, then the Mandarin for to arrange, then say \emph{zhàoxiàng} again
\end{itemize}

\begin{tcolorbox}[breakable,colback=teal!4,colframe=teal!35!black,title={Before you move on}]
What does zhàoxiàng mean? (\textquotedblleft{}To take a photo\textquotedblright{}.) Say it once more.
\end{tcolorbox}
\section[jǐngdiǎn]{jǐngdiǎn — sight; place of interest}

\label{lesson:ZH-C207-jingdian}



\begin{quote}
Before the new one: say the Mandarin for to arrange, then the Mandarin for to take a photo.
\end{quote}

\subsection*{You'll want to know: jǐngdiǎn}
\emph{jǐngdiǎn} — \textquotedblleft{}sight; place of interest\textquotedblright{}.

\begin{quote}
\emph{nǐ qùguò zhège jǐngdiǎn ma} — Have you been to this sight?
\end{quote}

Its characters come later; for now it is a word for the ear.

Two third tones meet in this word: jǐng is said with a rising tone before diǎn, and the spelling keeps both marks.

\begin{cousinweb}
\emph{jǐngdiǎn} is built from \emph{jǐng} 'a view' and \zh{点} 'a point', both inherited from Middle Chinese, \emph{kjaengX} and \emph{temX} (Guangyun, in Baxter's transcription), which glosses the first 'bright' and 'great'.
\end{cousinweb}

\begin{grammarlens}[title={What you've built}]
One more word for talking about plans.
\end{grammarlens}

\subsection*{Your turn}
\begin{itemize}
\raggedright
  \item \emph{Say it:} \emph{jǐngdiǎn}
  \item \emph{Say it:} \emph{jǐngdiǎn}, once more
  \item \emph{Say it:} \emph{nǐ qùguò zhège jǐngdiǎn ma} — Have you been to this sight?
  \item \emph{Recall:} say the Mandarin for to arrange, then the Mandarin for to take a photo, then say \emph{jǐngdiǎn} again
\end{itemize}

\begin{tcolorbox}[breakable,colback=teal!4,colframe=teal!35!black,title={Before you move on}]
What does jǐngdiǎn mean? (\textquotedblleft{}Sight; place of interest\textquotedblright{}.) Say it once more.
\end{tcolorbox}
